% =====================================================================
% Section: Derivation of the quantum postulates
% Drop-in for Paper 1. Requires in the preamble:
%   \usepackage{amsmath, amssymb}
%   \usepackage[backend=biber, style=phys]{biblatex}   % or your chosen style
% Internal cross-references (\ref{...}) are placeholders; repoint to your labels:
%     sec:complex-amplitudes -- the self-reference -> negative probability ->
%                               Shannon -> complex amplitudes derivation
%     sec:l2-metric-fibre    -- the L^2 / complex-Hopf section
%     sec:born-rule          -- the U(1)-invariance / minimum-degree section
%     sec:wigners-friend     -- the Wigner's-friend / FANOUT / redundancy section
%     sec:wf-reconciliation  -- the Tsirelson / PR-box aside
% Citation keys refer to the companion .bib; reconcile with your Mendeley keys.
% British English throughout.
% =====================================================================

\section{Derivation of the quantum postulates}
\label{sec:postulates}
We've now assembled enough pieces of the jigsaw to be able to derive the postulates of quantum mechanics

The Dirac--von Neumann postulates are ordinarily laid down as independent
assumptions. In this framework they are not assumptions but a single chain: each
follows from the Inside/Outside Equivalence Principle by one move from the
preceding step. We assemble the chain here and mark the status of each link as
proved, structural, or still owed, so that no step claims more than it has earned.

The chain begins below the postulates. An embedded observer, required by the
equivalence principle to describe a totality that includes itself, must assign a
probability to a proposition about its own future state. The Lawvere--Yanofsky
diagonal \autocite{Yanofsky2003} and Popper's predictor argument
\autocite{Popper1950a, Popper1950b} show that no consistent classical assignment
exists: the self-referential proposition is undecidable to the observer before
the fact. Representing the result of that undecidable self-prediction forces a
signed, non-classical probability, which carries the observer's irreducible
self-ignorance. From this single obstruction the postulates follow.

\paragraph{Postulate 1 (complex Hilbert space).}
Two independent steps fix the state space. The field is forced first: a signed
probability, read through the Shannon information measure and analytically
continued, yields a complex amplitude $S=-\ln\lvert p\rvert-\mathrm{i}\theta$,
whose phase $\theta$ is the self-referential ignorance
(Section~\ref{sec:complex-amplitudes}). The metric is forced second, and
separately: the principle of reversible continuity---that a continuous reversible
transformation connects any two pure states---selects the $L^{2}$ norm uniquely.
For every $\ell^{p}$ with $p\neq 2$ the linear isometries are generalised
permutations \autocite{Banach1932, Lamperti1958}, a group that cannot connect a
basis state to a superposition; only at $p=2$ is the norm-preserving group
connected and transitive on the sphere \autocite{Aaronson2004h}, and Hardy's
operational $K=N^{2}$ counting reaches the same quadratic norm by an independent
route \autocite{Hardy2001a}. The state manifold is therefore the unit sphere of the
Hermitian $2$-norm in $\mathbb{C}^{n}$, taken modulo the global phase
(Section~\ref{sec:l2-metric-fibre}). Placing the metric here, before the Born
rule rather than after it, is deliberate: the Born derivation presupposes the
round sphere, so the metric must be anchored independently to avoid circularity.
Two bounds belong with this step. Reversible continuity fixes the metric, not the
field---real and quaternionic quantum theory satisfy it equally---so the
restriction to $\mathbb{C}$ rests on the Shannon step, not on continuity. And
deriving reversible continuity itself from self-reference is the one arrow the
chain still owes; until it is closed, continuity enters as a principle, and the
weight of the $L^{2}$ conclusion is carried jointly with the independent
isometry and counting arguments.

\paragraph{Postulate 2 (Hermitian observables).}
A measurement result must be communicable: it must be copyable from one observer
to another, the \textsc{copy} operation of the diagrammatic calculus
\autocite{CoeckeKissinger2017}. The imaginary part of the amplitude is the
private self-referential phase and cannot be broadcast; only the real,
gauge-invariant content can. Observables whose spectra are real and therefore
communicable are exactly the self-adjoint operators. Hermiticity is thus the
condition that a quantity be the kind of thing a measurement can report.

\paragraph{Postulate 3 (Born rule).}
\label{sec:born-rule}
Outcome probabilities must be invariant under the global phase, the $U(1)$
self-ignorance, realised geometrically as the $S^{1}$ fibre of the complex Hopf
fibration $S^{1}\to S^{3}\to S^{2}$ over the Bloch sphere
\autocite{Hopf1931, Urbantke2003}. There is no $U(1)$-invariant of degree one; the
lowest-degree invariant is quadratic, so the probability must be
$p=\lvert\psi\rvert^{2}$. The squaring is forced
rather than merely permitted, and the argument holds for a single qubit, where
Gleason's theorem is silent. The Born quadratic is not a new structure: it is the
same $L^{2}$ form that fixed the metric in Postulate~1, now read as the
normalisation of the round sphere, and the fibre over which the phase is averaged
has constant length, so the phase is hidden uniformly
\autocite{GluckWarnerZiller1986}. This step is a structural argument rather than a
theorem in the strict sense, and is flagged as such.

Note: Using Hopf fibrations to get the Born rule is something Claude AI came up with on its own without my prompting. 

\paragraph{Postulate 4 (state reduction).}
Measurement is a \textsc{fanout}: the gauge-invariant relational content is copied
into a record, converting a private epistemological amplitude into a public
fact. For the observer who performs it, acquiring that fact is Bayesian updating,
and the projection postulate is its expression. The reduction is perspectival---a
fact relative to the observer in the sense of relational quantum mechanics
\autocite{Rovelli1996a}---and it is graded by redundancy. A single copy is a
relative fact, still reversible by an agent controlling the joint system; redundant
proliferation across many independent fragments makes the fact objective, which is
quantum Darwinism \autocite{zurek_2009} and the stable/relative-fact distinction of
Di~Biagio and Rovelli \autocite{DiBiagioRovelli2021}. \textquote{Collapse} and unitary
evolution are then the two charts---inside and outside the \textsc{fanout}
boundary---of one process, as the Wigner's-friend analysis shows
(Section~\ref{sec:wigners-friend}); the irreversibility and its $kT\ln 2$ cost
\autocite{Landauer_1961, wigner_1961} accrue with redundancy rather than at the
single act. The reduction postulate is therefore recovered as a chart-dependent
description, not a fundamental dynamical law; the identification of redundancy as
the grading parameter is a synthesis of the framework's parts and is flagged as
such.

\paragraph{Postulate 5 (unitary evolution).}
Conservation of probability is preservation of the $L^{2}$ inner product, and the
inner-product-preserving linear maps are the unitaries $U(n)$---the same
isometry fact that fixed the metric in Postulate~1. Continuity in time together
with Stone's theorem gives $U(t)=\mathrm{e}^{-\mathrm{i}Ht}$, the Schr\"odinger
equation. Reversibility is here a corollary, not a premise: bare reversibility on
the amplitude space yields only the general linear group, and it is preservation
of the quadratic form that carves $U(n)$ out of it.

\paragraph{Postulate 6 (composite systems).}
Systems combine by the tensor product because the division algebra norm is
multiplicative \autocite{hurwitz1898composition}: composition preserves the $L^{2}$ structure
only if states multiply. Cayley--Dickson doubling carries $\mathbb{C}$ to the next
rung, and the ascent of the Hilbert space is the ascent of the Hopf tower, complex
to quaternionic, used in the Wigner's-friend construction
\autocite{mosseri2001hopf}.

\paragraph{Consequence (the Tsirelson bound).}
The postulates fix not only the form of the theory but the strength of its
correlations. In the two-qubit ball the maximal entanglement sits at the centre,
and a Popescu--Rohrlich box \autocite{PopescuRohrlich1994} would require a point
beyond it; the CHSH correlator is capped at $2\sqrt{2}$
\autocite{Cirelson1980QuantumInequality, VerstraeteWolf2002, Held2021}, with the value recovered
at the centre and its general derivation still owed
(Section~\ref{sec:wf-reconciliation}).

\medskip
\noindent The result is that the six postulates are one chain, not six premises.
Of the links, the field and the Born quadratic rest on the framework's signed
probability and $U(1)$-invariance arguments; the $L^{2}$ metric, the unitary group
and the tensor product are over-determined by reversible continuity, the isometry
classification and Hurwitz composition; state reduction is recovered as the
chart-relative face of \textsc{fanout}. The two outstanding debts are the
derivation of reversible continuity from self-reference and the forcing of the
Tsirelson value as a theorem rather than a recovered number---named here so that a
reader can see exactly where the chain is complete and where it is not.